\section{Fixed-support degree criteria and punctured-observable obstructions}
\label{sec:nonlinear-fixed-support}
We use Boolean algebraic degree over $\mathbb F_2$, not real Walsh degree.
Let $S\subset\mathbb F_2^r$ have size $N$, let $T$ be its nonempty
complement, and suppose $N\equiv4\pmod8$, $r\ge4$. The seed law is
uniform on $S$. Put $t=\operatorname{affdim}(T)$ and $a=r-t$.
The low-degree model considered first is degree at most $r-2$.
The code background is standard Reed--Muller duality; see Carlet's
\emph{Boolean Functions for Cryptography and Coding Theory},
\url{https://www.math.univ-paris13.fr/~carlet/book-fcts-Bool-vect-crypt-codes.pdf}.
The duality and all required rank facts are proved below. No novelty
or new Hadamard order is asserted.

\begin{lemma}[Even restriction code and its radical]
\label{lem:nonlinear-radical}
Let $C=\mathrm{RM}(r-2,r)|_S$ be the binary restriction code, with
$C_0$ its even-weight subcode. Then
\[
 \dim C=N-a,\quad \dim C_0=N-a-1,\quad
 \dim\operatorname{rad}(C_0)=a+1,
\]
where the radical uses the binary dot product. Consequently its bilinear
rank is $N-2a-2$.
\end{lemma}
\begin{proof}
Monomials of combined degree less than $r$ have even full-cube sum;
dimension counting gives $\mathrm{RM}(r-2,r)^\perp=\mathrm{RM}(1,r)$.
Thus $C^\perp$ consists of the restrictions of affine polynomials
vanishing on $T$, a space of dimension $r-t=a$. Every such restriction
lies in $C$ and has even weight: its nonzero full-cube weight is $2^{r-1}$.
Also $\mathbf1\in C$ has even weight $N$.
The vector $\mathbf1$ is not in $C^\perp$, because that would make
$\mathbf1_S$ affine on the full cube, impossible at weight
$N\equiv4\pmod8$ with $r\ge4$. Hence the parity functional on $C$ is
nonzero, giving $\dim C_0=N-a-1$. Finally
$C_0^\perp=C^\perp+\langle\mathbf1\rangle\subseteq C_0$;
therefore the radical is this entire $(a+1)$-dimensional space.
\end{proof}

\begin{theorem}[Exact fixed-support criterion for a full basis]
\label{thm:nonlinear-fixed-support}
Suppose an order-$N$ Hadamard matrix exists. The support $S$ admits
$N-1$ unbiased pairwise-independent sign observables of degree at most
$r-2$ if and only if $T$ affinely spans $\mathbb F_2^r$.
If $T$ does span, every prescribed row labeling of any such Hadamard
matrix admits observables with this degree bound.
\end{theorem}
\begin{proof}
Let $B$ be the binary truth columns of the observables on $S$.
Balancedness gives weight $N/2$, even; pairwise orthogonality gives
intersection size $N/4$, odd. Hence $B^TB=I+J$ over $\mathbb F_2$.
For $N-1$ columns its rank is $N-2$. All columns lie in $C_0$, so
Lemma~\ref{lem:nonlinear-radical} forces
$N-2\le N-2a-2$, or $a=0$.
Conversely, $a=0$ means $C^\perp=0$, so restriction of
$\mathrm{RM}(r-2,r)$ to $S$ is surjective. Extend each labeled binary
Hadamard column separately.
\end{proof}

\paragraph{An explicit support-switching operation.}
This surjectivity can be implemented as a small syndrome repair, rather
than a new interpolation search. Relabel the existing rows arbitrarily
onto $S$, set provisional outside values to zero, and choose an affine
basis $z_0,\ldots,z_r$ inside $T$. Compute the $r+1$ binary parity defects
of the provisional truth table against $1,z_1,\ldots,z_r$ (coordinate
functions, with the basis points distinguished by their indices).
Flip values only at the chosen basis points to cancel those defects.
The matrix with columns $(1,z_i)$ is invertible, so there is a unique
repair on those points. The repaired function is orthogonal to all
affine functions and therefore has degree at most $r-2$.
No value on $S$ changes, so all real moments are preserved exactly.
For degree $r-1$, a single omitted point suffices: repair only total
truth-table parity. This operation works for any proper support and
permits arbitrary row relabeling.

\begin{corollary}[Degree saturation near a power of two]
\label{cor:nonlinear-degree-saturation}
Assume $2^{r-1}<N<2^r$, $N\equiv4\pmod8$, and $H_N$ exists.
There is some minimal-seed encoding of degree at most $r-2$ if and
only if $2^r-N\ge r+1$. If $2^r-N\le r$, the smallest possible maximum
degree is exactly $r-1$. In particular every existing
$H_{2^r-4}$, $r\ge4$, requires degree exactly $r-1$ on a minimal seed.
\end{corollary}
\begin{proof}
A set of at most $r$ omitted labels cannot affinely span an $r$-space;
a set of at least $r+1$ can be chosen to contain an affine basis.
Use the theorem, and the one-point parity repair for the upper bound.
\end{proof}

\subsection{The $N-4$-bit extraction input still requires growing degree}
\begin{theorem}[Punctured-observable bounds, without completion]
\label{thm:nonlinear-punctured}
Let $N\equiv4\pmod8$ and let $N$ equal-weight seed labels carry $m=N-4$
unbiased pairwise-independent sign bits, all of Boolean degree at most
$d$. Then
\[
 \sum_{j=0}^d\binom rj\ge N-3.
\]
At minimal $r=\lceil\log_2N\rceil$, this forces
$d\ge\lceil r/2\rceil$. For degree at most $r-2$, the sharper necessary
condition on the omitted labels is $\operatorname{affdim}(T)\ge r-1$.
For $N=2^r-4$ and $r\ge5$, degree $r-2$ is impossible for this punctured
input, not only for a full Hadamard basis.
\end{theorem}
\begin{proof}
The same binary Gram is $I+J_m$, now nonsingular because $m$ is even.
Thus the $m$ truth columns are independent. The constant vector is
orthogonal to them, is nonzero, and cannot belong to their nondegenerate
span. The Boolean restriction space must have dimension at least
$m+1=N-3$. Since $N-3>2^{r-1}$ at minimal seed length, binomial symmetry
proves the growing-degree bound.
For $d=r-2$, the nonsingular Gram embeds an $m$-dimensional nondegenerate
space into $C_0$. Its bilinear rank must therefore be at least $m$,
so $N-4\le N-2a-2$ implies $a\le1$. Four omitted points have affine
dimension at most three; when $r\ge5$, this violates $t\ge r-1$.
No completion theorem or existing full Hadamard matrix was used.
\end{proof}

\subsection{A sharp four-bit quadratic classification}
\begin{proposition}[Exact finite obstruction]\label{prop:nonlinear-twelve-classification}
For a uniform law on any twelve labels in $\mathbb F_2^4$, the largest
number of unbiased pairwise-independent quadratic sign observables is:
\[
 \begin{cases}
 5,&\text{if the four omitted labels form an affine plane},\\
 7,&\text{if they are affinely independent}.
 \end{cases}
\]
In particular the order-twelve $N-4=8$-bit extraction input requires
at least one cubic observable, for every labeling.
\end{proposition}
\begin{proof}[Exact finite certificate]
Affine input changes preserve degree and real moments. Four distinct
points in $\mathbb F_2^4$ have precisely two affine equivalence types:
representatives $\{0,1,2,3\}$ and $\{0,1,2,4\}$.
For each representative, enumerate the binary code generated by all
degree-at-most-two monomial restrictions to the twelve retained labels.
It has respectively $1024$ and $2048$ words. Balanced sign observables
are exactly weight-six words; identify a word with its complement by
requiring its first coordinate zero. This leaves respectively $126$
and $238$ vertices. Join two words when their xor has weight six.
Cliques are exactly sets of unbiased pairwise-independent observables.
The accompanying original standard-library checker exhaustively excludes
cliques of sizes six and eight, respectively, and verifies cliques of
sizes five and seven. Sorted-support integer-word witnesses are
\[
 \begin{array}{c|l}
 \{0,1,2,3\}&246,826,1372,1708,3178\\
 \{0,1,2,4\}&246,826,1372,1708,3178,3504,3606.
 \end{array}
\]
All enumeration and recursive candidate intersections are exact bitset
operations; the checker also recovers explicit quadratic polynomials.
This is an exhaustive finite computational certificate, not a numerical
optimization or a Lean formalization.
\end{proof}
Combining the proposition with Theorem~\ref{thm:nonlinear-punctured}
shows that the punctured input at every order $2^r-4$, $r\ge4$, requires
maximum degree at least $r-1$ in this equal-weight minimal-seed model.
Whenever a full Hadamard matrix at that order exists, one-point parity
repair gives matching degree $r-1$ examples. This obstruction does not
apply automatically to nonuniform measures, larger seeds, or approximate
moments; those are distinct possible search directions.
